\documentclass[10pt,twocolumn]{article}
\usepackage[margin=0.75in]{geometry}
\usepackage{graphicx}
\usepackage{booktabs}
\usepackage{amsmath}
\usepackage[hidelinks]{hyperref}
\usepackage{microtype}
\setlength{\parskip}{2pt}
\title{SLO-budgeted adaptive chunked-prefill sizing (single-node deflection analogue)}
\author{a02 (auto-inference)}
\date{2026-09-03}
\begin{document}
\maketitle
\begin{abstract}
The mechanism sizes each prefill chunk against the live decode
time-between-tokens slack rather than a static token count, on the premise that
this workload's TTFT budget (p90 $\le$ 2818\,ms) is generous enough to convert
into deeper decode batches.  Attempt~1 first measured the cost curve the
mechanism depends on and found it flat --- 60.2 / 61.6 / 61.7\,$\mu$s per
prefilled token at chunk 2048 / 4096 / 8192, a 2.5\% spread over a $4\times$
range --- so what ships is an SLO-constrained sizer over a three-arm ladder that
settles on the configured size and is thereafter identical to stock here.  It
prices at \$6.9392/1k at $N^{*}{=}12$ against the \$8.32/1k baseline
($-16.6$\%) with all accuracy gates passing, but an inert diff cannot earn that,
so the delta is not attributable to the mechanism.  Chunk-size policy cannot
move cost per output token on this stack, because prefill cost per token barely
depends on chunk size.
\end{abstract}

\section{Mechanism}
At fixed GPU price, cost per output token is $\mathrm{mean\_TPOT}/B$, and
\[
\mathrm{mean\_TPOT} = a_d + b_d B + \frac{P\,a_e + L/R}{\mathrm{steps/cycle}} .
\]
The decode step is a fixed cost plus a per-sequence KV read; the shipped priors
are $a_d = 10.5$\,ms and $b_d = 0.9$\,ms/seq, both refit online from the
scheduler's loop period.  $L/R$ and $B$ are set by the trace and the SLO.  The
only term a chunk-sizing policy owns is $P$, the number of extend steps, at
$a_e$ each --- so the idea reduces to: how large is $a_e$, and does ms per
prefilled token vary with chunk size?

Both were measured.  On workbench-5 a fresh 20{,}480-token prompt prefilled as
$8192{+}8192{+}4096$ took 1287.0\,ms end to end while the three chunk forwards
themselves sum to 1285.1\,ms, so $a_e \approx 2$\,ms --- not the $\sim$10\,ms
weight re-read the linear picture suggests.  \texttt{ncu} says why: the prefill
GEMM runs at 82.9\% of SM peak against 28.1\% of DRAM peak, so it is compute
bound and the 23.8\,GB weight read is already amortised inside the GEMM.  The
intercept a linear fit reports is not a cost you can stop paying.

\section{What changed}
Attempt~0 built the idea as recorded --- chunk solver \emph{plus} deliberate
admission delay --- and failed the SLO gate without pricing.  Two deviations
follow, both deliberate.  Admission delay is a pure loss here: the market
workload is a closed loop of $N$ users, so holding a request back cannot densify
a later batch (the user is blocked, nothing arrives to fill the gap) and
concurrency falls 1:1 with throughput; the original TTFT win comes from
cross-node queueing and KV transfer, both zero on one H100, so that half is not
built.  And sizing chunks against \emph{instantaneous} slack shrinks them
exactly when the batch is busy: prefill work is invariant to chunk size, so each
split adds a step without removing work, raising mean TPOT on the tightest
batches.

What ships replaces the constant at the one point it enters the scheduler,
\texttt{PrefillAdder(...,\,chunked\_prefill\_size,\,...)}: a ladder
\texttt{[static/4, static/2, static]}, page-aligned and snapped to a 512-token
grid so the stack sees a few recurring prefill shapes rather than a new one each
step, with \texttt{largest\_arm\_within\_budget} taking the largest arm whose
predicted stall fits the tail-TBT credit the tightest in-flight decode has
banked.  Growth past the configured size sits behind
\texttt{--adaptive-chunk-max-growth} (default 1, shrink only): raising
\texttt{chunked\_prefill\_size} to 24576 at \texttt{mem\_fraction\_static} 0.84
OOMs inside \texttt{PrefillCudaGraphRunner.capture()} (workbench-6), since
prefill graphs are captured at shapes derived from it.  Per-request TBT
bookkeeping hangs off \texttt{Req}, is maintained by the scheduler loop rather
than the request path, and is cleared on retract.  Four files under
\texttt{srt/} change; no launch-line change is recorded.

A search variant was built and cut.  It ran a bandit over the ladder to
\emph{measure} the curve, and it produced the numbers above: on the live server
at stock defaults it reported \texttt{2048:60.2us/tok(n=11)},
\texttt{4096:61.6(n=46)}, \texttt{8192:61.7(n=91)} over 148 extend steps, then
latched after $\sim$53 prefill passes having spent $\sim$24 off-static ones.
The equivalence gate rejected it at decode agreement 0.7744 against stock's
1.0000.  \texttt{rem\_chunk\_tokens} is not only the per-request chunk cap but
also the budget the adder packs \emph{other} waiting requests into, so probing
it changes prefill batch composition, and FP8 block GEMMs are not
batch-invariant.  Scoping the probe to continuation passes made it worse, not
better ($0.8128 \rightarrow 0.7744$): with 32 prompts submitted at once, stock
itself chunks one, so the probe still fired.  Paying numerics drift to search a
curve already known to be flat is a bad trade.

\section{Results}
\begin{table}[h]\centering\small
\begin{tabular}{llrrrl}
\toprule
attempt & tier & \$/1k & $\Delta$\% & N* & gates \\
\midrule
0 & screen & - & - & None & slo \\
1 & screen & 6.73 & -25.4 & 12 & gsm8k 64\%, longbench 49\%, mmlu 66\% \\
1 & full & 6.94 & -16.6 & 12 & gsm8k 64\%, longbench 49\%, mmlu 66\% \\
1 & full & 6.94 & -16.6 & 12 & gsm8k 64\%, longbench 49\%, mmlu 66\% \\
\bottomrule
\end{tabular}
\caption{Priced attempts. Baseline \$8.32/1k.}
\end{table}
\begin{figure}[h]\centering\includegraphics[width=\linewidth]{fig0.png}\caption{attempt-001: slo-mean-tpot}\end{figure}
\begin{figure}[h]\centering\includegraphics[width=\linewidth]{fig1.png}\caption{attempt-001: price-vs-share}\end{figure}
\begin{figure}[h]\centering\includegraphics[width=\linewidth]{fig2.png}\caption{attempt-001: slo-p90-ttft}\end{figure}

Attempt~0 never priced.  Attempt~1 held $N^{*}=12$ at
\$6.9392440795727115/1k full-tier and \$6.734769414252982/1k screen-tier
($-25.42$\% against the screen tier's own baseline, which the record does not
state).  The full row appears twice with every digit identical: one sweep
recorded twice, not a repeat.

The figures confirm the premise and refute the lever.  Fig.~3: p90 TTFT never
approaches its 2818\,ms limit at any level swept --- the largest observed is
under 900\,ms --- and throughput is uncorrelated with it ($r^2 = 0.009$), so the
slack the idea wanted to spend is real.  Fig.~1: mean TPOT binds instead and
sets $N^{*}$; $N{=}12$ sits at 16.4\,ms against the 20\,ms limit, $N{=}16$ at
22.1\,ms.  Fig.~2: price is exactly inverse in served share ($r^2 = 1.000$,
\$5.15/1k per point of share), so cost is a pure function of the $N$ the SLO
admits, and adjacent levels of that quantised grid are about 15\% apart
($\approx$\$8.4 at $N{=}8$, \$6.94 at 12, $\approx$\$5.8 at 16, read off
Fig.~2).

That is the difficulty in one line.  Prefill is 3--14\% of GPU-seconds here and
chunk size moves prefill by at most 2.5\%, so the ceiling on the mechanism is
$0.14 \times 0.025 = 0.35$\% of cost per output token --- under single-sweep
noise, and two orders of magnitude below one grid step.  Since the shipped sizer
settles on the configured size, the priced $-16.6$\% measures the gap between
this sweep's stock behaviour and the recorded \$8.32/1k fleet baseline, not the
change.

\section{What the gates said}
All accuracy gates passed on both tiers: GSM8K 64\%, LongBench F1 49\%, MMLU
66\%.  No decode-agreement or $|\Delta\log p|$ figure for the shipped variant
appears in the record available here; the only agreement numbers measured on
this idea are the cut search variant's, 0.8128 unscoped and 0.7744 scoped,
against stock's 1.0000.

Exercise matters more than the pass.  The shipped sizer settles on the
configured chunk size and is thereafter identical to stock on this trace, so the
gates drove no non-stock path: a pass here is evidence the change is inert, not
that it is correct under load.  The one variant that did schedule chunks
differently is exactly the one equivalence caught.

\section{Next}
The experiment this argues for is not another chunk-sizing variant but a clean
re-price of stock on this exact stack and sweep.  A diff that never leaves the
stock chunk size returned $-16.6$\% against the \$8.32/1k reference, so that
reference is stale by roughly that much and every result priced against it
inherits the error.  Cost is one priced sweep at the same five levels
($N \in \{4,8,12,16,24\}$); the record gives no separate dollar or wall-clock
figure for a sweep.

That re-priced reference should then be aimed at the decode step, not the
scheduler.  Moving $N^{*}$ from 12 to 16 needs mean TPOT at $N{=}16$ to fall
from 22.1\,ms below 20\,ms, a 9.5\% cut at fixed batch, and with
$a_d = 10.5$\,ms and $b_d = 0.9$\,ms/seq that is a decode-kernel or KV-bytes
question.  The TTFT slack is genuine; nothing in prefill scheduling converts it.

\section*{Erratum: ablation by the harness operator (2026-09-03)}
\emph{Added after the agent finished; not the agent's text.} The conclusion
above, that the shipped sizer is inert here and the $-16.6\%$ therefore
indicts a stale baseline, was tested directly and is wrong on both counts.
\begin{table}[h]\centering\small
\begin{tabular}{lrrr}
\toprule
stack at $N=12$, 120 s & \$/1k & mean TPOT & fwd s \\
\midrule
stock (baseline-v3b) & 8.32 & 19.8 ms & 122.7 \\
this diff, sizer on (attempt 1) & 6.94 & 16.4 ms & 109.7 \\
this diff, sizer on (replicate) & 6.82 & 16.4 ms & 110.0 \\
this diff, \texttt{SGLANG\_DISABLE\_ADAPTIVE\_CHUNK=1} & 8.30 & 19.4 ms & 122.6 \\
\bottomrule
\end{tabular}
\caption{Same four files, same launch line; only the kill switch differs
(\texttt{runs/a02-ablate-kill}). The sizer is the entire effect.}
\end{table}
The baseline is not stale: the diff with its sizer disabled reproduces it
to \$0.02. The sizer is not a no-op at the deployed point: with it on, the
same 123 s wall serves 7--9\% more output tokens in 13 fewer seconds of
forward time, mean TPOT falls from 19.8 to 16.4 ms at the same batch
($\approx$11.5--12.0) and hit rate, and decode forward time per output token
falls from 2.17 to 1.78 ms. The credit projection walks the ladder down
at $N=12$; why a scheduler-only change moves decode forward time per token
by 18\% is the open question this result actually poses.

On the gates: the equivalence check was run on the shipped stack
(\texttt{runs/a02-b4-equiv}, 32 LongBench prompts, 1216 positions):
top-1 agreement 1.0000, $|\Delta\log p|$ 0.0000, and greedy 64-token
decode agreement 1.0000 with every prompt exact. The shipped variant is
token-equivalent to stock and is exercised under load; the pass is
evidence of correctness, not of inertness.

\end{document}
